%--------------------------
\begin{frame}{Question 1}

\begin{block}{}
La demande internationale cesse de produire ses effets en phase nationale si :

\begin{itemize}
  \item[$\circ$] a. le délai de 30 mois à compter de la date de priorité n’est pas respecté pour 
  l’ouverture de la phase nationale. 
  \item[$\circ$] b. la taxe nationale n’est pas payée dans le délai de 31 mois à compter de la date de 
  priorité. 
  \item[$\circ$] c. les signatures manquantes ne sont pas remises dans le délai de 30 ou 31 mois à 
  compter de la date de priorité. 
  \item[$\bullet$] d. Aucune des réponses a, b ou c n’est correcte.
\end{itemize}

\end{block}

\end{frame}

%--------------------------
\begin{frame}{Question 1}

\begin{itemize}
  \item[$\circ$] \textbf{a.} \; Faux : le délai d’entrée en phase nationale n’est pas toujours
de 30 mois, certains États, LU et TZ, appliquant encore un délai de 20 mois
(\pctart{22}.1)).

  \item[$\circ$] \textbf{b.} \; Faux : certains États, LU et TZ, appliquant encore un délai de 20 mois
  (\pctart{22}.1)).

  \item[$\circ$] \textbf{c.} \; Faux : la signature de tout
  déposant qui n'a pas signé la requête est une condition
  permise par le PCT mais pas obligatoire
  (\pctart{27}.2) + \pctrule{51bis}.1.a.i)) 

  \item[$\circ$] \textbf{d.} \; Correct puisqu'aucune des réponses \textbf{a},
  \textbf{b} ou \textbf{c} est correcte.
\end{itemize}

\end{frame}
